% GENERATED FILE. Edit canonical lessons, then run npm run generate:books.
% canonical-source-hash: fnv1a64:323f727049d308f5
% canonical-lessons: KA-C86-oneminute, KA-C86-wait, KA-C86-show, KA-C86-didnthear, KA-C86-didntunderstand

\chapter{One Minute, Please}
\label{ch:ka-one-minute-please}
\hlchaptermodality{86}

\begin{chapteropening}
\textbf{By the end of this chapter:} \emph{I can ask for a minute, ask someone to wait or show me, and say I did not hear or did not understand.}

Complete the last lesson of chapter 86: i can ask for a minute, ask someone to wait or show me, and say I did not hear or did not understand.
\end{chapteropening}

\section[ondu nimiṣa]{\kn{ಒಂದು} \kn{ನಿಮಿಷ} (ondu nimiṣa) — one minute}

\label{lesson:KA-C86-oneminute}



\begin{quote}
Before the new one: say the Kannada for meaning, then the Kannada for please write it.
\end{quote}

\subsection*{You'll want to know: \kn{ಒಂದು} \kn{ನಿಮಿಷ}}
\textbf{\kn{ಒಂದು} \kn{ನಿಮಿಷ}} (\emph{ondu nimiṣa}) — \textquotedblleft{}one minute\textquotedblright{}. \textbf{\kn{ಒಂದು}} is \textquotedblleft{}one\textquotedblright{}, from the numbers chapter; \textbf{\kn{ನಿಮಿಷ}} (\emph{nimiṣa}) is \textquotedblleft{}minute\textquotedblright{}.

Say it with a raised hand when you need a moment.

\begin{grammarlens}[title={What you've built}]
A number you know and a new word for time.
\end{grammarlens}

\subsection*{Your turn}
\begin{itemize}
\raggedright
  \item \emph{Say it:} \emph{ondu nimiṣa}
  \item \emph{Say it:} \emph{ondu nimiṣa}, once more
  \item \emph{Say it:} say \emph{ondu nimiṣa}
  \item \emph{Recall:} say the Kannada for meaning, then the Kannada for please write it, then say \emph{ondu nimiṣa} again
\end{itemize}

\begin{tcolorbox}[breakable,colback=teal!4,colframe=teal!35!black,title={Before you move on}]
What does \kn{ಒಂದು} \kn{ನಿಮಿಷ} mean? (\textquotedblleft{}One minute\textquotedblright{}.) Say it once more.
\end{tcolorbox}
\section[kāyiri]{\kn{ಕಾಯಿರಿ} (kāyiri) — please wait}

\label{lesson:KA-C86-wait}



\begin{quote}
Before the new one: say the Kannada for please write it, then the Kannada for one minute.
\end{quote}

\subsection*{You'll want to know: \kn{ಕಾಯಿರಿ}}
\textbf{\kn{ಕಾಯಿರಿ}} (\emph{kāyiri}) — \textquotedblleft{}please wait\textquotedblright{}, the polite form of \textbf{\kn{ಕಾಯು}} (\emph{kāyu}), \textquotedblleft{}to wait\textquotedblright{}.

\textbf{\kn{ಒಂದು} \kn{ನಿಮಿಷ} \kn{ಕಾಯಿರಿ}} (\emph{ondu nimiṣa kāyiri}) — \textquotedblleft{}please wait a minute\textquotedblright{}.

\begin{grammarlens}[title={What you've built}]
The polite ending ‑\kn{ಇರಿ} asks gently.
\end{grammarlens}

\subsection*{Your turn}
\begin{itemize}
\raggedright
  \item \emph{Say it:} \emph{kāyiri}
  \item \emph{Say it:} \emph{kāyiri}, once more
  \item \emph{Say it:} say \emph{ondu nimiṣa kāyiri}
  \item \emph{Recall:} say the Kannada for please write it, then the Kannada for one minute, then say \emph{kāyiri} again
\end{itemize}

\begin{tcolorbox}[breakable,colback=teal!4,colframe=teal!35!black,title={Before you move on}]
What does \kn{ಕಾಯಿರಿ} mean? (\textquotedblleft{}Please wait\textquotedblright{}.) Say it once more.
\end{tcolorbox}
\section[tōrisi]{\kn{ತೋರಿಸಿ} (tōrisi) — please show me}

\label{lesson:KA-C86-show}



\begin{quote}
Before the new one: say the Kannada for one minute, then the Kannada for please wait.
\end{quote}

\subsection*{You'll want to know: \kn{ತೋರಿಸಿ}}
\textbf{\kn{ತೋರಿಸಿ}} (\emph{tōrisi}) — \textquotedblleft{}please show (me)\textquotedblright{}.

\textbf{\kn{ದಾರಿ} \kn{ತೋರಿಸಿ}} (\emph{dāri tōrisi}) — \textquotedblleft{}please show me the way\textquotedblright{}, with \textbf{\kn{ದಾರಿ}}, \textquotedblleft{}way, road\textquotedblright{}.

\begin{grammarlens}[title={What you've built}]
When words fail, ask to be shown.
\end{grammarlens}

\subsection*{Your turn}
\begin{itemize}
\raggedright
  \item \emph{Say it:} \emph{tōrisi}
  \item \emph{Say it:} \emph{tōrisi}, once more
  \item \emph{Say it:} say \emph{dāri tōrisi}
  \item \emph{Recall:} say the Kannada for one minute, then the Kannada for please wait, then say \emph{tōrisi} again
\end{itemize}

\begin{tcolorbox}[breakable,colback=teal!4,colframe=teal!35!black,title={Before you move on}]
What does \kn{ತೋರಿಸಿ} mean? (\textquotedblleft{}Please show me\textquotedblright{}.) Say it once more.
\end{tcolorbox}
\section[kēḷisalilla]{\kn{ಕೇಳಿಸಲಿಲ್ಲ} (kēḷisalilla) — I didn't hear}

\label{lesson:KA-C86-didnthear}



\begin{quote}
Before the new one: say the Kannada for please wait, then the Kannada for please show.
\end{quote}

\subsection*{You'll want to know: \kn{ಕೇಳಿಸಲಿಲ್ಲ}}
\textbf{\kn{ಕೇಳಿಸಲಿಲ್ಲ}} (\emph{kēḷisalilla}) — \textquotedblleft{}(it) was not heard\textquotedblright{}, the way to say \textquotedblleft{}I didn't hear\textquotedblright{}. It is built on \textbf{\kn{ಕೇಳು}}, \textquotedblleft{}hear, listen\textquotedblright{}, with \textbf{‑\kn{ಇಲ್ಲ}}, \textquotedblleft{}not\textquotedblright{}.

\begin{grammarlens}[title={What you've built}]
A known verb with a known no.
\end{grammarlens}

\subsection*{Your turn}
\begin{itemize}
\raggedright
  \item \emph{Say it:} \emph{kēḷisalilla}
  \item \emph{Say it:} \emph{kēḷisalilla}, once more
  \item \emph{Say it:} say \emph{kēḷisalilla, mattomme hēḷi}
  \item \emph{Recall:} say the Kannada for please wait, then the Kannada for please show, then say \emph{kēḷisalilla} again
\end{itemize}

\begin{tcolorbox}[breakable,colback=teal!4,colframe=teal!35!black,title={Before you move on}]
What does \kn{ಕೇಳಿಸಲಿಲ್ಲ} mean? (\textquotedblleft{}I didn't hear\textquotedblright{}.) Say it once more.
\end{tcolorbox}
\section[arthavāgalilla]{\kn{ಅರ್ಥವಾಗಲಿಲ್ಲ} (arthavāgalilla) — I didn't understand}

\label{lesson:KA-C86-didntunderstand}



\begin{quote}
Before the new one: say the Kannada for please show, then the Kannada for I didn't hear.
\end{quote}

\subsection*{You'll want to know: \kn{ಅರ್ಥವಾಗಲಿಲ್ಲ}}
\textbf{\kn{ಅರ್ಥವಾಗಲಿಲ್ಲ}} (\emph{arthavāgalilla}) — \textquotedblleft{}I didn't understand\textquotedblright{}, literally \textquotedblleft{}the meaning did not happen\textquotedblright{}. It starts with \textbf{\kn{ಅರ್ಥ}}, \textquotedblleft{}meaning\textquotedblright{}, from the last chapter.

\begin{grammarlens}[title={What you've built}]
Meaning, and the no that says it did not arrive.
\end{grammarlens}

\subsection*{Your turn}
\begin{itemize}
\raggedright
  \item \emph{Say it:} \emph{arthavāgalilla}
  \item \emph{Say it:} \emph{arthavāgalilla}, once more
  \item \emph{Say it:} say \emph{arthavāgalilla, nidhānavāgi hēḷi}
  \item \emph{Recall:} say the Kannada for please show, then the Kannada for I didn't hear, then say \emph{arthavāgalilla} again
\end{itemize}

\begin{tcolorbox}[breakable,colback=teal!4,colframe=teal!35!black,title={Before you move on}]
What does \kn{ಅರ್ಥವಾಗಲಿಲ್ಲ} mean? (\textquotedblleft{}I didn't understand\textquotedblright{}.) Say it once more.
\end{tcolorbox}
